\documentclass[pmlr,twocolumn,10pt]{jmlr} % ML4H 2026 Proceedings track (PMLR/JMLR W&CP style)

\usepackage{caption}
% Packages auto-loaded by jmlr: amsmath, amssymb, natbib, graphicx, url, algorithm2e
\usepackage{booktabs}
\usepackage{float}
\usepackage{amsmath}
\usepackage{siunitx}

% Required for line numbers during review; remove \linenumbers before camera-ready.
\usepackage[switch]{lineno}

\jmlrvolume{XXX} % PMLR volume number, assigned upon acceptance
\jmlryear{2026}
\jmlrworkshop{Machine Learning for Health (ML4H) 2026}

\title[Feature-Set Scaling for Kidney Failure Prediction]{Feature-Set Scaling for Kidney Failure Prediction: Benchmarking Survival Models Against the Kidney Failure Risk Equation on MIMIC-IV}

% Anonymized for double-blind review per ML4H 2026 policy.
\author{Anonymous Author(s)}

\linenumbers % Activate line numbering for review; remove for camera-ready.

\begin{document}

\maketitle

\begin{abstract}
The Kidney Failure Risk Equation (KFRE) uses four or eight clinical variables to estimate kidney failure risk. We benchmark ten learned survival models and the published KFRE equation on MIMIC-IV-derived cohorts with four features (2{,}809 patients), eight features (1{,}213 patients), and twenty laboratory features with missingness indicators (32{,}601 patients; no KFRE analogue). Five runs per scenario provide 160 model--scenario--run results. Patient-level discrimination is summarized as mean (sample SD); the runs reuse a fixed patient partition, so SD measures fitting and tuning variation rather than patient-sampling uncertainty. Dynamic-DeepHit has the highest recorded C-index: 0.827 (0.011), 0.880 (0.009), and 0.856 (0.035), respectively. KFRE has C-index 0.539 in both applicable scenarios. Hazard Transformer reaches 0.717 (0.060) on the full twenty-feature cohort, superseding its poor feasibility-pilot result. These rankings do not imply uniformly good native risk calibration: Dynamic-DeepHit substantially underpredicts risk in one twenty-feature run, and Hazard Transformer yields constant two-year probabilities. Differences in cohorts, representations, and available patient histories prevent isolation of feature-count effects. Below-chance predictions, test-set reuse for Logistic Hazard tuning, and row-based censoring references remain evaluation limitations. Decision-curve superiority claims are withheld because model and baseline curves use different evaluation units. The completed runs provide a descriptive benchmark and identify the validation changes required before conclusions about clinical utility can be drawn.
\end{abstract}

\begin{keywords}
Survival analysis, chronic kidney disease, Kidney Failure Risk Equation, MIMIC-IV, benchmark
\end{keywords}

\paragraph*{Data and Code Availability}
This work uses MIMIC-IV \citep{johnson2023mimic_ehr}, available via PhysioNet to credentialed users under its data use agreement. Anonymized code implementing the cohort extraction, models, and analyses is provided as supplementary material for review; a public repository will be released upon publication.

\paragraph*{Institutional Review Board (IRB)}
No IRB approval was required. MIMIC-IV is a de-identified, publicly available critical care database released under a data use agreement; this study involved no additional identifiable patient contact.

\section{Introduction}\label{sec:intro}

Chronic kidney disease (CKD) affects over 1 in 7 American adults and can progress to end-stage renal disease (ESRD) \citep{CDC2023CKD}. Early identification of high-risk patients is critical for nephrology referral and dialysis-access planning \citep{levin2011early_detection}. The clinical standard for this task is the Kidney Failure Risk Equation (KFRE) \citep{tangri2011predicting}, a closed-form Cox model built from as few as four routinely collected variables (age, sex, eGFR, urine albumin-to-creatinine ratio (uACR)), or eight in its extended form (adding calcium, phosphate, bicarbonate, albumin). KFRE has been externally validated across 31 cohorts spanning over 700{,}000 patients \citep{tangri2016multinational} and is embedded in clinical guidelines \citep{major2019kfre_external_validation}.

Modern EHR resources and flexible survival models motivate comparison with KFRE across different input representations. Adding laboratory variables can also change cohort eligibility and the available history, complicating interpretation of feature-count comparisons.

We benchmark ten learned survival models in three MIMIC-IV-derived scenarios: \textbf{four\_features} and \textbf{eight\_features} match KFRE's four- and eight-variable inputs; \textbf{twenty\_features} uses 20 laboratory variables with missingness flags and has no published KFRE analogue. We add the closed-form KFRE equation to the first two scenarios (11 models there). The extracted cohorts contain 2{,}809, 1{,}213, and 32{,}601 patients, respectively. They share a source population but do not form a controlled, nested feature-count experiment.

Building complete KFRE-input rows requires merging asynchronous measurements around creatinine anchors while retaining repeated observations (Section~\ref{sec:cohort}). Nearest-value matching can include measurements after an anchor, so this design remains retrospective.

\textbf{Contributions.} We provide a documented cohort-construction procedure and a completed comparison across all three scenarios, reporting mean and sample SD over five runs on the same patient partition. We use one prediction per patient while distinguishing final-row and full-history inputs. We also document the limits of the completed analysis: scenario-specific cohorts, test-informed LogisticHazard selection, raw-row censoring references, and unmatched decision-curve comparators. The results describe these implementations and do not establish the isolated effect of feature count or clinical utility.

\section{Methods}\label{sec:methods}

\textbf{Problem formulation.} We frame progression to ESRD as right-censored time-to-event estimation. Each covariate row $\boldsymbol{x}_i\in\mathbb{R}^d$ is paired with a duration $t_i\in\mathbb{R}^+$ and event indicator $e_i\in\{0,1\}$. Patient histories are stored in a counting-process (\texttt{start}, \texttt{stop}) format, supporting time-varying Cox fitting, snapshot-based KFRE prediction, and sequence models consuming a patient's full row history (Dynamic-DeepHit, Hazard Transformer).

\subsection{Cohort} \label{sec:cohort}
The source population is 34{,}332 MIMIC-IV patients meeting a CKD stage 3--5 or ESRD diagnosis-code criterion (57.2\% male, mean age 67.8, SD 16.1).

\textbf{Four/eight\_features (KFRE-comparable).} KFRE's inputs come from separate, asynchronously-ordered labs (creatinine, urine albumin/creatinine, and for the 8-variable form, a chemistry panel), so assembling complete covariate rows required a merge design (full derivation and literature citations in Appendix~\ref{app:cohort}): each row anchors on one creatinine draw (eGFR computed at that instant via CKD-EPI 2021); calcium/phosphate/bicarbonate/albumin are matched to the nearest value within $\pm$24h (following APACHE II's precedent for a concurrent chemistry-panel snapshot); uACR is matched to the nearest value across the patient's whole history (loosened from a same-admission bound after that bound caused 96.5\% cohort attrition with outcome-rate skew); rows missing a required lab are dropped, not imputed. This keeps four\_features/eight\_features complete-case and time-varying (one row per qualifying creatinine draw), with scenario-specific cohort selection. Matching is retrospective and can include measurements after the creatinine anchor.

\textbf{Twenty\_features (heterogeneous).} We select the 20 most frequently measured labs for this cohort, each represented as a (value, missingness-flag) pair (40 columns), one row per lab event, no cross-lab merge — a broader sparse-input scenario that also changes cohort eligibility and representation, rather than isolating feature count.

Table~\ref{tab:cohort_flow} reports the distinct extracted cohorts. The eight-feature cohort (1{,}213 patients, 1{,}032 events) is the smallest; these full-cohort counts do not establish training-sample adequacy for all architectures.

\begin{table}[t]
\centering
\scriptsize
\begin{tabular}{@{}lcccc@{}}
\toprule
& \textbf{Source} & \textbf{four} & \textbf{eight} & \textbf{twenty} \\
\midrule
Patients & 34,332 & 2,809 & 1,213 & 32,601 \\
Records & 74,197 & 26,829 & 5,407 & 8.13M \\
\% male & 57.2 & 58.3 & 60.4 & 57.1 \\
Age, yrs (SD) & 67.8 (16.1) & 63.9 (15.1) & 62.1 (15.4) & 67.8 (16.1) \\
ESRD, \% & 91.9 & 83.5 & 85.1 & 91.5 \\
Train/test pts & --- & 2247/562 & 970/243 & 26080/6521 \\
\bottomrule
\end{tabular}
\caption{Cohort flow: source population vs.\ each scenario's final cohort.}
\label{tab:cohort_flow}
\end{table}

\subsection{KFRE Baseline} \label{sec:kfre}
For four/eight\_features we add KFRE as a non-trained benchmark: $\mathrm{Risk}(t)=1-S_0(t)^{\exp(L)}$, $S_0$(2yr)$=0.9750$ (four variables) or $0.9780$ (eight variables) \citep{tangri2016multinational}, with
\begin{align}
L_4 &= -0.2201(\tfrac{\text{age}}{10}{-}7.036) + 0.2467(\text{male}{-}0.5642) \nonumber \\
&\ \ -0.5567(\tfrac{\text{eGFR}}{5}{-}7.222) + 0.4510(\ln(\text{uACR}){-}5.137)
\end{align}
for the 4-variable form and an analogous 8-term $L_8$ (adding albumin, phosphate, bicarbonate, calcium terms) for the 8-variable form \citep{tangri2011predicting} (full coefficients in Appendix~\ref{app:kfre}).

\subsection{Benchmarked Models} \label{sec:models}
We benchmark 10 learned models per scenario, plus KFRE on four/eight\_features only (11 total there): \textbf{Cox} (\texttt{lifelines}' \texttt{CoxTimeVaryingFitter}, fit directly on the counting-process rows) and \textbf{Weibull AFT} \citep{cox1972regression,anderson1991nonproportional_weibull,davidson2019lifelines}; the neural models \textbf{DeepSurv} \citep{katzman2018deepsurv}, \textbf{Dynamic-DeepHit} \citep{lee2020dynamicdeephit,lee2018deephit}, \textbf{RNN-Surv} (LSTM-based) \citep{giunchiglia2018rnn_surv}, \textbf{LogisticHazard} (\texttt{pycox}'s off-the-shelf discrete-time hazard model), and \textbf{Hazard Transformer} \citep{hazard_transformer,vaswani2017attention}; and the ensembles \textbf{GBSA} \citep{chen2013gradient_GBSA}, \textbf{Random Survival Forest} \citep{pickett2021random_survival_forests,ishwaran2008random}, and \textbf{FastSurvivalSVM}, via \texttt{scikit-survival} \citep{polsterl2020sksurv}. Every architecture's exact structural form, loss, and tuned hyperparameter range — verified directly against the model implementations, not restated from a fixed reference configuration — is given in Appendix~\ref{app:models}; neural models use 10-trial Optuna searches, GBSA/SRF use grid search, and FastSurvivalSVM/Cox/Weibull AFT use fixed hyperparameters. Neural selection is based on training C-index except LogisticHazard, which uses the designated test partition for validation and trial/checkpoint selection; its test result is therefore subject to model-selection leakage (Appendix~\ref{app:models}).

\subsection{Evaluation} \label{sec:eval}
All three scenarios are reported across five runs on the same 80/20 patient-ID partition (\texttt{random\_state=42}), as mean $\pm$ sample SD ($n=5$). These SDs quantify run-to-run variation, not independent-cohort uncertainty. Every model yields one prediction per patient, but snapshot models use final rows while Dynamic-DeepHit and Hazard Transformer use full histories. Follow-up remains anchored to the first lab rather than a prospective prediction landmark.

Harrell's C-index uses full available follow-up. Mean cumulative/dynamic AUC evaluates model-specific scalar scores daily over days 1--728; IBS integrates over 50 times from days 1--730. IBS uses a training-fitted, shifted-score Breslow-style survival conversion for every model, including KFRE, rather than native probabilities. IBS/AUC censoring weights use raw training rows, a remaining mismatch with patient-level test predictions (Appendix~\ref{app:evaluation}).

Calibration at two and five years uses native probabilities when supported and the score conversion otherwise, with Kaplan--Meier risk per decile. Generated decision curves \citep{vickers2006decisioncurve} compare patient-level model outputs with raw-row treat-all/eGFR baselines and inconsistent eGFR decision thresholds; they cannot support clinical-utility comparisons. Competing mortality was not modeled. Full definitions and limitations appear in Appendix~\ref{app:evaluation}.

\section{Results}\label{sec:results}

All 160 model--scenario--run metric records are available across five runs: 11 models each for four/eight features and 10 for twenty features. Tables~\ref{tab:discrimination} and~\ref{tab:twenty} report mean (sample SD) on the fixed patient partition. Rep99 pilot results are excluded; SD measures run variation, not uncertainty across independent test cohorts.

\subsection{Four and Eight Features}
\begin{table}[t]
\centering
\scriptsize
\setlength{\tabcolsep}{3pt}
\begin{tabular}{@{}lccc@{}}
\toprule
\multicolumn{4}{c}{\textbf{Four features}} \\
\textbf{Model} & \textbf{C-index}$\uparrow$ & \textbf{IBS}$\downarrow$ & \textbf{Mean AUC}$\uparrow$ \\
\midrule
Cox & 0.308 (0.000) & 0.531 (0.000) & 0.272 (0.000) \\
Dynamic-DeepHit & \textbf{0.827 (0.011)} & \textbf{0.279 (0.018)} & \textbf{0.939 (0.017)} \\
Hazard Transformer & 0.637 (0.001) & 0.409 (0.007) & 0.701 (0.002) \\
Logistic Hazard & 0.551 (0.012) & 0.416 (0.006) & 0.567 (0.018) \\
RNN-Surv & 0.522 (0.006) & 0.430 (0.011) & 0.526 (0.011) \\
DeepSurv & 0.540 (0.020) & 0.463 (0.075) & 0.549 (0.023) \\
GBSA & 0.518 (0.003) & 0.438 (0.008) & 0.530 (0.008) \\
Survival RF & 0.532 (0.005) & 0.445 (0.003) & 0.543 (0.007) \\
Survival SVM & 0.543 (0.000) & 0.417 (0.000) & 0.571 (0.000) \\
Weibull AFT & 0.552 (0.000) & 0.414 (0.000) & 0.579 (0.000) \\
KFRE & 0.539 (0.000) & 0.552 (0.000) & 0.553 (0.000) \\
\midrule
\multicolumn{4}{c}{\textbf{Eight features}} \\
\textbf{Model} & \textbf{C-index}$\uparrow$ & \textbf{IBS}$\downarrow$ & \textbf{Mean AUC}$\uparrow$ \\
\midrule
Cox & 0.470 (0.000) & 0.578 (0.000) & 0.425 (0.000) \\
Dynamic-DeepHit & \textbf{0.880 (0.009)} & \textbf{0.281 (0.024)} & \textbf{0.953 (0.009)} \\
Hazard Transformer & 0.605 (0.039) & 0.529 (0.205) & 0.631 (0.049) \\
Logistic Hazard & 0.521 (0.019) & 0.490 (0.016) & 0.500 (0.029) \\
RNN-Surv & 0.531 (0.005) & 0.493 (0.005) & 0.590 (0.015) \\
DeepSurv & 0.484 (0.014) & 0.555 (0.115) & 0.486 (0.022) \\
GBSA & 0.541 (0.011) & 0.529 (0.014) & 0.559 (0.011) \\
Survival RF & 0.522 (0.008) & 0.523 (0.007) & 0.539 (0.008) \\
Survival SVM & 0.540 (0.000) & 0.488 (0.000) & 0.595 (0.000) \\
Weibull AFT & 0.554 (0.000) & 0.485 (0.000) & 0.597 (0.000) \\
KFRE & 0.539 (0.000) & 0.737 (0.000) & 0.579 (0.000) \\
\bottomrule
\end{tabular}
\caption{Mean (sample SD), five runs on a fixed split. IBS uses the common scalar-risk transformation. C-index spans follow-up; IBS/AUC summarize the two-year interval. Bold: best recorded mean. SD is run variability, not a confidence interval.}
\label{tab:discrimination}
\end{table}

Dynamic-DeepHit has the highest recorded C-index, 0.827 (0.011) and 0.880 (0.009), and mean AUC, 0.939 (0.017) and 0.953 (0.009), for four and eight features. KFRE has C-index 0.539 in both scenarios. Hazard Transformer follows at 0.637 (0.001) and 0.605 (0.039). The remaining non-Cox models cluster at 0.518--0.552 and 0.484--0.554. These descriptive comparisons do not establish equivalence or statistical superiority.

Four-to-eight-feature changes are model-dependent: Dynamic-DeepHit +0.053, Cox +0.162, DeepSurv -0.056, and Hazard Transformer -0.032 in mean C-index; the other seven models change by 0.03 or less. The cohorts differ, so this is not an isolated test of adding four measurements. Cox remains below 0.5 (0.308/0.470); DeepSurv's eight-feature mean is also below 0.5.

\subsection{Full Twenty-Feature Cohort}
\begin{table}[t]
\centering
\scriptsize
\setlength{\tabcolsep}{3pt}
\begin{tabular}{@{}lccc@{}}
\toprule
\multicolumn{4}{c}{\textbf{Twenty features (heterogeneous)}} \\
\textbf{Model} & \textbf{C-index}$\uparrow$ & \textbf{IBS}$\downarrow$ & \textbf{Mean AUC}$\uparrow$ \\
\midrule
Cox & 0.457 (0.000) & 0.397 (0.000) & 0.433 (0.000) \\
Dynamic-DeepHit & \textbf{0.856 (0.035)} & \textbf{0.248 (0.013)} & \textbf{0.946 (0.026)} \\
Hazard Transformer & 0.717 (0.060) & 0.314 (0.032) & 0.802 (0.078) \\
Logistic Hazard & 0.522 (0.005) & 0.374 (0.006) & 0.531 (0.026) \\
RNN-Surv & 0.523 (0.001) & 0.376 (0.003) & 0.527 (0.001) \\
DeepSurv & 0.533 (0.003) & 0.370 (0.001) & 0.539 (0.007) \\
GBSA & 0.521 (0.000) & 0.371 (0.000) & 0.529 (0.000) \\
Survival RF & 0.522 (0.001) & 0.372 (0.000) & 0.534 (0.001) \\
Survival SVM & 0.513 (0.000) & 0.369 (0.000) & 0.540 (0.000) \\
Weibull AFT & 0.525 (0.000) & 0.370 (0.000) & 0.546 (0.000) \\
\bottomrule
\end{tabular}
\caption{Full twenty-feature cohort: mean (sample SD), five runs, 6{,}521 test patients. Metric definitions follow Table~\ref{tab:discrimination}.}
\label{tab:twenty}
\end{table}

The full cohort contains 32{,}601 patients and approximately 8.13 million lab-event rows. Dynamic-DeepHit leads at C-index 0.856 (0.035), mean AUC 0.946 (0.026), and transformed-score IBS 0.248 (0.013). Hazard Transformer reaches C-index 0.717 (0.060), replacing the earlier poor pilot result. The other non-Cox models span 0.513--0.533; Cox is 0.457. Cohort and representation differences prevent attribution to feature count alone.

\subsection{Importance and Calibration Diagnostics}
Figure~\ref{fig:feature_importance_four} illustrates rep1 importance. eGFR ranks first for four neural models, whereas GBSA emphasizes uACR and eGFR similarly. Raw coefficients, permutation scores, and gradients use different scales; a uACR coefficient rounded to 0.0000 does not imply that the input was discarded. All five twenty-feature DDH/Transformer gradient analyses failed memory allocation; their fallback values are excluded from interpretation.

\begin{figure}[t]
\centering
\includegraphics[width=\linewidth]{figs/four_features_all_models_feature_importance.png}
\caption{Four-feature importance diagnostics, rep1. Model-specific measures are not comparable across panels.}
\label{fig:feature_importance_four}
\end{figure}

Native risk calibration varies across runs. For twenty-feature rep1, Dynamic-DeepHit's 2-year risk-group means are 0.166--0.187 versus observed risks of 0.431--0.998, despite strong discrimination. Hazard Transformer also produces constant or near-constant horizon probabilities. The reported IBS uses transformed scalar scores and does not evaluate these native probabilities directly.

Decision-curve comparisons are withheld: model curves use patients, whereas treat-all/eGFR baselines use raw rows; the eGFR rule also uses a different threshold definition. The saved curves therefore do not establish clinical utility.

\section{Discussion}\label{sec:discussion}

Dynamic-DeepHit has the highest recorded mean discrimination across all three scenarios, followed by Hazard Transformer. The completed twenty-feature results replace the pilot findings. However, feature count varies with cohort composition, missingness representation, and access to patient histories; these experiments cannot isolate the benefit of additional measurements or attribute performance differences to architecture alone.

Native calibration is less consistent than discrimination, and the common transformed-score IBS does not measure native probability calibration. The decision-curve comparisons require harmonized patient-level baselines before any referral benefit can be inferred.

\subsection{Limitations}\label{sec:limitations}
Five runs on a fixed partition describe fitting/tuning variability, not patient-sampling uncertainty; there are no patient-bootstrap confidence intervals or external validation. The event-rich hospital cohort limits generalizability. IPCW uses raw training rows, and Logistic Hazard uses the nominal test set in tuning/checkpoint selection. These evaluation choices require validation in a rerun with separate tuning data and patient-level censoring estimates.

Models also differ in access to longitudinal histories. Final-observation scoring, whole-history uACR matching, and bidirectional chemistry windows do not establish prospective performance at a common decision time. The uACR match attaches a value recorded after the anchor creatinine for 40\% of rows (median absolute separation 264 days, p95 5.4 years), whereas the $\pm$24h chemistry matches are simultaneous with the anchor for 93--99\% of rows. For label-positive patients the median time from the first extracted record to the labelled event is 1.2 days (four-feature), 0.4 days (eight-feature) and 104 days (twenty-feature), and 18--35\% contribute a single row, so a two-year horizon measured from that origin is a follow-up window rather than a prediction lead time. Cox remains below chance across scenarios, eight-feature DeepSurv averages below 0.5, and Hazard Transformer has constant two-year probabilities; shared implementation effects remain possible. The labelled outcome is an acute/unspecified kidney-failure code, not the dialysis-or-transplant endpoint KFRE predicts (Section~\ref{sec:outcome}), so KFRE-versus-model calibration differences combine an endpoint difference, a population difference, and the fitted/externally-specified asymmetry, which this benchmark does not separate. Death before ESRD is not modeled as a competing event, demographic subgroup performance is untested, and twenty-feature DDH/Transformer importance outputs are failed-probe fallbacks. Future work should address these issues before clinical utility is assessed.


\bibliography{sn-bibliography}

\appendix

\section{Cohort Construction: Full Derivation} \label{app:cohort}

KFRE's inputs are drawn from separate, asynchronously-ordered labs (serum creatinine, a urine albumin/creatinine test, and, for the 8-variable form, a chemistry panel) — unlike a design that avoids multi-lab merging altogether by encoding one row per single-lab event with missingness flags (the twenty\_features scenario's approach, Section~\ref{sec:cohort}). Producing complete covariate rows from asynchronously measured inputs for four\_features/eight\_features required a merge design with two design choices: (1) whether the row format should still be time-varying given that KFRE itself is a point-in-time equation, and (2) how to handle a patient having multiple draws of the same lab within one admission.

\textbf{Row format.} We keep the time-varying (\texttt{start}, \texttt{stop}) format used throughout. KFRE's formula is agnostic to row packaging — it can be applied independently to each row, treating every row as its own landmark snapshot in the sense of landmarking methodology for time-varying covariates in survival analysis \citep{vanhouwelingen2007landmarking,putter2022landmarking2}.

\textbf{Merge rule: creatinine-anchored nearest-value matching.} Each row is anchored on one creatinine measurement (eGFR is computed from it, via the CKD-EPI 2021 equation, at that instant). Other required labs are attached to the anchor by nearest-value matching using the matching rules below rather than by unbounded last-observation-carried-forward, which is known to bias sparse EHR-derived time-varying covariates \citep{hu2022locf_bias}:
\begin{itemize}
    \item \textit{Calcium, phosphate, bicarbonate, serum albumin} (8-variable scenario only, routinely drawn on the same chemistry panel as creatinine): nearest value within $\pm$24 hours of the anchor. This window follows APACHE II's precedent of treating all labs drawn within a patient's first 24 hours — including serum creatinine — as one physiologic snapshot \citep{knaus1985apacheii}; no directly analogous 24h-snapshot convention was found for other candidate windows (e.g., KDIGO's 48h AKI change-detection window measures a \emph{trend}, not concurrency, and is not equivalent).
    \item \textit{uACR} (a separately-ordered urine test): nearest value across the patient's entire history. We initially bounded this match to the same hospital admission, following the resolution Tangri et al.\ document for one of their own real-world validation cohorts (\textit{"Geisinger: covariates obtained most closely to index date within a past year were included in models"} \citep{tangri2016multinational}), but this bound caused 96.5\% patient attrition (34{,}332 $\to$ 1{,}220 for the 4-feature scenario) and skewed the outcome rate from 91.9\% to 98.2\% ESRD-positive by disproportionately dropping ESRD-negative patients. We loosened the bound to the patient's whole history, increasing cohort size at the cost of temporal alignment. Because matching is not restricted to measurements before the anchor, these retrospective inputs can include future information relative to that anchor.
    \item Any row missing a required lab within its window is \emph{dropped}, not imputed or flagged — four\_features/eight\_features are complete-case by construction, matching KFRE's own derivation cohort's eligibility criteria \citep{tangri2011predicting}.
\end{itemize}
One row is emitted per qualifying creatinine draw, so a patient monitored multiple times within one admission correctly yields multiple rows, each independently paired with its own nearest in-window value of every other required lab (a sparser lab such as uACR is therefore legitimately reused, at its nearest value, across several rows for the same patient — expected under the landmarking framework above, not a bug).

Table~\ref{tab:cohort_flow} follows the two-column cohort-flow reporting convention of Major et al.\ \citep{major2019kfre_external_validation} and STROBE item 13(a) \citep{vonelm2007strobe}. All three scenarios are event-rich by construction (83--92\% patient-level ESRD-positive), because the source population is itself drawn from a CKD/ESRD diagnosis-code criterion.

\textbf{Cohort size.} The eight-feature scenario is the smallest, with 1{,}213 patients and 1{,}032 observed events. These counts describe the full extracted cohort, before its training/test split, and do not establish the adequacy of the training sample for every architecture. Complete-case selection also changes cohort composition, so cross-scenario differences cannot be attributed to feature count alone.

\section{KFRE Coefficients: Full 8-Variable Form} \label{app:kfre}
\begin{align}
L_8 &= -0.1992\left(\tfrac{\text{age}}{10}-7.036\right) + 0.1602(\text{male}-0.5642) \nonumber \\
&\quad -0.4919\left(\tfrac{\text{eGFR}}{5}-7.222\right) + 0.3364\left(\ln(\text{uACR})-5.137\right) \nonumber \\
&\quad -0.3441(\text{albumin}-3.997) + 0.2604(\text{phosphate}-3.916) \nonumber \\
&\quad -0.07354(\text{bicarb.}-25.57) -0.2228(\text{calcium}-9.355)
\end{align}
Coefficients from \citet{tangri2011predicting}; original-equation baseline survival constants from \citet{tangri2016multinational}: $S_0$(2yr)$=0.9780$ and $S_0$(5yr)$=0.9301$ for eight variables, compared with $0.9750$ and $0.9240$ for four variables. KFRE has no published equation for the twenty\_features scenario and is therefore absent from it.

\section{Benchmarked Models: Full Architectural Detail} \label{app:models}
The neural models, GBSA, and SRF are tuned per scenario (below), so we describe each model's structural form and loss precisely, verified directly against the model implementations, and give the tuned hyperparameter range rather than a single fixed value.

\subsection{Traditional Survival Analysis Models}
Both use \texttt{lifelines} \citep{davidson2019lifelines}; Cox is fitted to counting-process rows and Weibull AFT to each patient's final row.

\textbf{Cox proportional hazards} \citep{cox1972regression} models the hazard of subject $i$ at time $t$ as $\lambda_i(t) = \lambda_0(t)\exp(\boldsymbol{x}_i(t)^\intercal\boldsymbol{\beta})$, where $\boldsymbol{x}_i(t)$ is subject $i$'s (possibly time-varying) covariate row and $\lambda_0(t)$ an unspecified baseline hazard, fit by maximizing the Efron-corrected partial likelihood over all subjects' \texttt{start}/\texttt{stop} intervals directly (\texttt{lifelines}' \texttt{CoxTimeVaryingFitter}, \texttt{id\_col=subject\_id}) rather than one static row per patient — the model consumes every landmark row (Appendix~\ref{app:cohort}), not just each patient's last. We use an L2 penalizer of 1.0 to stabilize the fit against the strongly collinear rows a single patient contributes.

\textbf{Weibull Accelerated Failure Time (AFT)} \citep{anderson1991nonproportional_weibull} instead parameterizes the full survival distribution, $S(t\mid\boldsymbol{x}) = \exp(-(t/\lambda(\boldsymbol{x}))^{\rho})$, $\lambda(\boldsymbol{x}) = \exp(\boldsymbol{x}^\intercal\boldsymbol{\beta})$, fit on the flattened one-row-per-patient (last-observation) frame via \texttt{lifelines}' \texttt{WeibullAFTFitter}. For twenty\_features (whose 40, mostly-zero-filled columns produce a near-singular design matrix), we add an L2 penalizer of 0.1 on the scale parameter's coefficients, the library's own documented remedy; four\_features/eight\_features use the unpenalized default.

\subsection{Deep Survival Analysis Models}
DeepSurv and LogisticHazard use each patient's final row. Dynamic-DeepHit and Hazard Transformer consume full patient histories. RNN-Surv is trained on individual rows represented as length-one sequences and is evaluated on each patient's final row; its recurrent architecture therefore does not model longitudinal histories in this implementation.

\textbf{DeepSurv} \citep{katzman2018deepsurv} is a feedforward network — 1 to 20 \texttt{Linear}$\to$\texttt{ReLU}$\to$\texttt{Dropout} blocks (width 16--256, per-layer dropout 0.1--0.5, all tuned) followed by a single linear output — trained with the Cox partial log-likelihood on its scalar output, so its raw output is already a "higher = riskier" hazard score by construction.

\textbf{Dynamic-DeepHit} \citep{lee2020dynamicdeephit,lee2018deephit} feeds each patient's full covariate row sequence with a padding mask directly into a 2-layer \emph{bidirectional} LSTM (hidden width 16--128, dropout 0.1--0.5, tuned). The concatenated forward/backward outputs at each timestep pass through a \texttt{Linear}$\to$\texttt{Tanh} compression, then a masked additive (Bahdanau-style) attention layer pools the sequence into one context vector per patient, letting the model weight informative visits over uninformative ones. The pooled vector feeds 0--3 further cause-specific \texttt{Linear}$\to$\texttt{ReLU}$\to$\texttt{Dropout} layers (width 16--128, tuned) and a final per-risk linear head producing a probability mass function (PMF) over 5{,}475 daily bins (15 years). Following DeepHit's original formulation \citep{lee2018deephit}, training minimizes a composite loss $L=\alpha L_1+(1-\alpha)L_2$: $L_1$ is the negative log-likelihood of the discretized, right-censored outcome, and $L_2$ is a pairwise ranking loss rewarding the model for assigning an event subject higher cumulative incidence, at that subject's own event time, than any subject not yet failed by then — with $\alpha\in[0.1,1.0]$ tuned per scenario.

\textbf{RNN-Surv} \citep{giunchiglia2018rnn_surv} embeds each timestep through 1--3 linear embedding layers with intervening \texttt{ReLU} activations (width 32--128, tuned) into a 1--3-layer \emph{LSTM} (hidden width 64--256, tuned). A final linear layer maps the last hidden state to logits over $K\in[10,50]$ (tuned) discrete intervals spanning the maximum training follow-up, softmax-normalized into a PMF. As with Dynamic-DeepHit, the loss combines a discrete-time likelihood (observed-interval event probability for events; survival-beyond-interval probability for censored subjects) with a pairwise ranking term over the cumulative incidence function, weighted by a tunable coefficient in $[0.1,0.9]$.

\textbf{LogisticHazard} discretizes training follow-up into 50 bins (\texttt{pycox}'s \texttt{LabTransDiscreteTime}) and learns the log-odds of failure in each bin, via \texttt{pycox}'s off-the-shelf \texttt{MLPVanilla} network (1--3 hidden layers, width $\in\{16,32,64,128\}$, dropout 0--0.5, all tuned, batch size $\in\{32,64,128,256\}$) wrapping its \texttt{LogisticHazard} discrete-time hazard head and cross-entropy-style loss.

\textbf{Hazard Transformer} \citep{hazard_transformer,vaswani2017attention} departs from a sequence-attention reading of "Transformer": each patient's rows are linearly embedded and \emph{mean-pooled} (masked over observed rows) into a single patient representation \emph{before} any attention is applied. This pooled vector is broadcast across 100 fixed query time points spanning $[0,730]$ days, each combined with a learned time-encoding and a sinusoidal positional encoding; 2--6 (tuned) Transformer encoder layers (1--8 heads, tuned) then self-attend \emph{across these 100 time queries}, not across the patient's visits. A per-risk decoder produces a softmax-normalized PMF over the 100 bins (an earlier per-bin independent-sigmoid parametrization collapsed to a constant prediction and was replaced by this bounded-PMF form, the same fix Dynamic-DeepHit's risk heads use). Training minimizes the discrete-time event/censoring negative log-likelihood without an additional ranking term.

\subsection{Ensemble-Learning-Based Models}
All three are \texttt{scikit-survival} \citep{polsterl2020sksurv} estimators fit on the flattened last-observation frame.

\textbf{Gradient Boosting Survival Analysis (GBSA)} \citep{chen2013gradient_GBSA} fits an additive ensemble of regression trees against the Cox partial-likelihood gradient; grid-searched over \texttt{n\_estimators}$\in\{50,100,200,300\}$, \texttt{max\_depth}$\in\{3,5,10,15\}$, \texttt{learning\_rate}$\in\{0.01,0.1,0.2\}$ by cross-validated C-index.

\textbf{Random Survival Forest (SRF)} \citep{pickett2021random_survival_forests,ishwaran2008random} grows an ensemble of survival trees splitting on the log-rank statistic; grid-searched over \texttt{n\_estimators}$\in\{50,100,150\}$, \texttt{max\_depth}$\in\{\text{None},5,10\}$. Since \texttt{scikit-survival} does not expose impurity-based feature importance for this estimator (known to be biased for survival forests), we compute permutation importance for it directly.

\textbf{FastSurvivalSVM} learns a linear ranking function over the covariate space, optimizing a convex upper bound on the concordance index subject to margin and $\ell_2$ constraints. Fit with \texttt{scikit-survival}'s defaults ($\alpha=0.01$), as with the fixed-parameter Cox and Weibull AFT baselines.

\subsection{Model Training and Optimization}
The neural models use 10-trial Optuna searches per scenario, maximizing the implemented C-index objective. DeepSurv, Dynamic-DeepHit, RNN-Surv, and Hazard Transformer select trials on training-set C-index and stop when training loss fails to improve (patience five epochs, maximum 50 epochs). LogisticHazard is fitted for 100 epochs and uses the designated test partition both as \texttt{val\_data} and for trial/checkpoint selection. Its reported test performance therefore has model-selection leakage and is descriptive rather than an unbiased held-out estimate. GBSA and SRF use five-fold grid search within the training partition, selected by C-index. FastSurvivalSVM, Cox, and Weibull AFT use fixed hyperparameters. All neural models use Adam.

\section{Evaluation Definitions and Scope} \label{app:evaluation}
\textbf{Repetitions and aggregation.} We report all three scenarios from the completed clinical-validity analyses for repetitions 1--5. Each scenario uses an 80/20 split by patient ID with \texttt{random\_state=42}; the repetitions reuse the same partition. Results are the arithmetic mean and sample standard deviation ($n=5$, denominator $n-1$) of the five report values. They describe run-to-run variation on a fixed partition, not uncertainty across independent cohorts or externally validated confidence intervals.

\textbf{Prediction unit.} Every model contributes one test prediction per patient. Snapshot models use the final available row; Dynamic-DeepHit and Hazard Transformer use the available row history and the outcome recorded on its final row. RNN-Surv uses the final row as a length-one sequence. Thus the scoring unit is shared, while input histories differ. Durations remain measured from the extraction's first-lab anchor rather than being reset at a prospective prediction landmark. This retrospective evaluation does not establish prospective two-year prediction from a common baseline.

\textbf{Metrics.} Harrell's C-index is calculated over each model's available patient follow-up, without truncation at 730 days, after orienting scalar scores so higher values indicate greater risk. Mean cumulative/dynamic AUC uses that scalar score on a daily grid from day 1 through day 728 (or the shorter observed support), with the survival-weighted mean returned by \texttt{scikit-survival}. Score definitions remain model-specific: for example, Dynamic-DeepHit uses day-730 cumulative incidence, Hazard Transformer uses day-365 cumulative incidence, RNN-Surv uses expected event earliness, and LogisticHazard evaluates survival near the scored partition's median duration.

The integrated Brier score (IBS) uses 50 equally spaced times from day 1 to day 730. For every model, including KFRE, the reported IBS uses a score-derived survival curve $\hat S(t\mid r)=\exp[-r^*\hat H_0(t)]$, where $r^*=\max(r-\min(r_{\mathrm{train}}),10^{-6})$ and $\hat H_0$ is a Breslow-style cumulative baseline hazard fitted to that model's training-patient scores and outcomes. These IBS values evaluate this additional conversion rather than each model's native survival probabilities. For IBS only, test follow-up beyond the maximum training duration is administratively censored there. The inverse-probability-of-censoring reference for IBS and AUC is constructed from raw training rows, while test predictions are per patient; this mismatch remains a limitation of the completed analysis.

\textbf{Calibration and exploratory decision curves.} Calibration compares mean predicted risk with Kaplan--Meier-observed risk within risk deciles at 730 and 1{,}825 days. It uses native horizon probabilities where supported, with the score-derived survival conversion above as fallback. Unlike the IBS comparison, native KFRE calibration therefore uses the published equation directly. The generated decision-curve analyses \citep{vickers2006decisioncurve} include treat-all, treat-none, and eGFR referral rules ($<$30 or $<$45), but model curves use patient-level predictions whereas treat-all and eGFR comparators use raw test rows. The eGFR calculation also uses its referral fraction as an implied risk threshold. These unmatched units and thresholds preclude clinical-utility comparisons; the generated curves are exploratory and are not used to infer superiority over referral baselines. A competing-risk analysis for death before ESRD was not performed because exported durations omit the absolute calendar anchor needed to align death dates.


\end{document}
